\section{Referencial teórico}
\begin{enumerate}
    \item\label{rt_migrations} Migrações - as migrações são classes que representam interações pontuais com a base de dados, geralmente utilizadas para versionamento da estrutura da base de dados ou para alterações pontuais nos dados quando necessário.
    \item\label{rt_orm} Mapeador objeto-relacional (ORM) - os mapeadores objeto-relacional são abstrações que tem como premissa servir de ponte para duas linguagens, no caso do Vortex seriam o PHP que é uma linguagem orientada a objetos e o SQL que é a linguagem da base de dados. O mapeador age traduzindo o contexto de um lado da ponte para outro e vice-versa.
    \item\label{rt_reflections} \textit{Reflections} -   definida por \cite{gabriel1991clos}, reflexão é a capacidade de um programa manipular como dados algo que representa o estado do programa durante sua própria execução. Há dois elementos nessa manipulação: introspecção e intercessão. A introspecção refere-se à capacidade de um programa observar e, consequentemente, raciocinar sobre seu próprio estado. A intercessão é a habilidade de um programa modificar seu próprio estado de execução ou alterar sua própria interpretação ou significado. Ambos os aspectos requerem um mecanismo para codificar o estado de execução como dados; esse processo é chamado de reificação.
    \item\label{rt_framework} \textit{Framework} - conjunto de abstrações e implementações que visam facilitar ou incrementar a produção de determinada aplicação.
    \item\label{rt_guessers} \textit{Guessers} - são objetos ou métodos relacionados à ''adivinhação'' ou estimativas com base em lógicas predefinidas.
\end{enumerate}